\section{Evaluation criteria}
\label{sec:background}

\subsection{Mean Dice}

For validation record $i$, $\Omega_i=\{1,\ldots,H_i\}\times
\{1,\ldots,W_i\}$ is the set of pixel coordinates of an image with height $H_i$ and
width $W_i$, that is, its discrete domain. Let $P_i\subseteq\Omega_i$ be the predicted
foreground mask and $G_i\subseteq\Omega_i$ the union of its annotated instances. The
Dice coefficient is
\begin{equation}
D_i=\frac{2|P_i\cap G_i|+\varepsilon}
{|P_i|+|G_i|+\varepsilon},
\end{equation}
where $\varepsilon$ avoids division by zero. If both masks are empty, the implementation
assigns $D_i=1$. Over the $N$ annotator-image records in the validation set, the reported
metric is
\begin{equation}
\mathrm{Mean\ Dice}=\frac{1}{N}\sum_{i=1}^{N}D_i.
\end{equation}
This union-mask measure is appropriate for the strong class imbalance because it does
not reward correctly classified background pixels. However, it does not describe how
the foreground is divided into separate objects: a filament split into several
components can still retain substantial pixel overlap.

\subsection{Panoptic Quality}

PQ is therefore used as a complementary instance-level metric. For every
annotator-image record, predicted and annotated instances are associated when their
Intersection over Union is greater than 0.5. With this threshold, valid associations are
one-to-one. True matches and unmatched predicted or annotated instances are accumulated
over the validation split as TP, FP and FN. The metric is
\begin{equation}
\mathrm{PQ}=
\underbrace{\frac{\sum_{(p,g)\in\mathrm{TP}}\mathrm{IoU}(p,g)}{|\mathrm{TP}|}}_{\mathrm{SQ}}
\times
\underbrace{\frac{|\mathrm{TP}|}{|\mathrm{TP}|+\tfrac12|\mathrm{FP}|+\tfrac12|\mathrm{FN}|}}_{\mathrm{RQ}}.
\end{equation}
Segmentation Quality (SQ) measures the mean overlap of matched instances, whereas
Recognition Quality (RQ) penalises missing and spurious instances. The distinction is
relevant in this study because configurations with similar Mean Dice can produce
different numbers of fragmented or false-positive components.

The two measures are reported for one explicitly defined operating point at a time. In
particular, the best Dice from one output configuration is never combined with the best
PQ from another.
